% Claims that depend on the moved-goal, greedy and learning outcomes (written after the registered analyses completed).
\newcommand{\MGabstract}{When the goal moves, frozen human-choice agents still loop less but also head back to the old goal: their gain shrinks to \PMvLegal\ pp.}
\newcommand{\MGresults}{All frozen agents now succeed below a uniform policy (\MVhumanbc\% for human choices, \MVhumanlegaluniform\% legal, \MVcoldppo\% scratch, \MVratbc\% rat choices), because they carry the old goal with them: human-choice agents enter (3,3) in \MVoldgoalH\ of trials versus \MVoldgoalL\ for legal labels (Fig.~\ref{fig:diag}c, Fig.~\ref{fig:showcase}d). They still beat legal labels by \PMvLegal\ pp ($[\PMvLegallo,\PMvLegalhi]$, Holm $p=\PMvLegalp$) and scratch by \PMvScratch\ pp, and this residual gain is entirely their lower trapping (\MVtrapD\ pp fewer trapped trials; zero gain at equal trapping). Rat choices are slightly worse than legal labels (\PMvRat\ pp). Frozen agents therefore carry two things from human choices: how to move, which transfers, and where the goal was, which does not.}
\newcommand{\GRresults}{The human advantage holds under greedy actions too (\PGreedy\ pp, $[\PGreedylo,\PGreedyhi]$).}
\newcommand{\MGdiscussion}{With a moved goal the same habit still helps, but the remembered goal pulls frozen agents back, so new goals need learning (Sec.~\ref{sec:learning}).}
\newcommand{\LRresults}{All learners acquire the new goal (Fig.~\ref{fig:learning}); PPO ends at 54--67\% success and the searched local rules at 40--47\% from legal-label or scratch starts. The human-choice prior never speeds this learning. Under PPO it is unresolved ($-7.2$ pp learning-curve area versus legal labels, $[-19.5,5.1]$); under local rules that change only the top of the network it slows learning (TD-readout $-6.9$ pp, $[-11.8,-1.9]$; searched readout $-9.4$ pp; searched last-layer random feedback $-9.8$ pp; Holm $p\leq.04$), and the all-layer error-routing rule shows no difference ($-2.5$ pp, $[-8.9,3.9]$). The inherited pull toward the old goal is hardest to overwrite when only the readout can change. The local rules learn from the reward contingency, not from extra experience: permuting rewards among the parallel rooms removes TD-readout learning for every prior (final 23--39\% $\to$ 9--17\%, all $p<10^{-4}$), the counterpart of the permuted-decoding control for living neural cultures~\cite{zhou2026embodied}.}
\newcommand{\KPIresults}{\textbf{Forgetting.} The two goals share views, so nothing in a view says which goal is active, and every sequential learner overwrites part of the old one (Fig.~\ref{fig:kpi}a). Human-choice agents keep the most (post hoc, versus legal labels: PPO $+13.0$ pp, $[1.1,24.9]$; TD-readout $+13.9$, $[5.5,22.2]$). Trained on both goals at once, the same networks keep 79--88\% (diamonds): the loss is interference, not capacity. Locality alone does not help: local rules with replay keep 3\%. Brain mechanisms that supply the missing context do (c). Consolidation with old replay~\cite{benna2016,mcclelland1995} learns $+17$ pp more than PPO ($[9,25]$) at equal retention, and human choices still add $+10.0$ pp of retention (post hoc, $[4.0,16.0]$). A latent-cause gate~\cite{masse2018xdg,gershman2010latent} that receives no cue keeps \GTold\% versus 4.7\% ungated ($[\GTlo,\GThi]$ pp) and learns the new goal in 99.7\% of runs. A mushroom-body-like circuit~\cite{litwinkumar2017} with only three-factor plasticity and awake replay~\cite{foster2006reverse} learns the new goal better than PPO ($+12.4$ pp, $[3.8,21.1]$; registered searches, one test) while keeping as much of the old (TOST $p=.02$); with the gate it keeps $+13.3$ pp more of the old goal ($[6.9,19.7]$). \textbf{Generalisation.} The belief is allocentric, so learning in one room carries to the paired unseen room (3\% before; 9--34\% PPO and TD rules; 40--47\% with the expansion; b). \textbf{Spiking.} \SNNresults}
\newcommand{\SNNresults}{Converted to integrate-and-fire networks, human-choice agents keep their advantage (48\% versus 29\% legal at 8 steps; $+8.4$ pp, $[4.0,12.8]$, at 2 steps, with 15\% fewer operations than the dense network; d). Stopping each decision at a confidence bound, as crow and primate choices stop at an evidence bound~\cite{schmidbauer2026crow}, human-choice agents decide sooner ($-0.5$ steps, $[-0.7,-0.2]$, registered) at equal success. Learned networks convert without loss (crosses). Counts are operations, not energy.}
\newcommand{\GTold}{43}\newcommand{\GTlo}{+33.5}\newcommand{\GThi}{+42.4}
